\documentclass[11pt,letterpaper]{article}

\usepackage[spanish,es-noshorthands]{babel}
\usepackage{fontspec}
\setmainfont{Source Serif Pro}
\usepackage{microtype}
\usepackage{amsmath}
\usepackage{geometry}
\usepackage{booktabs}
\usepackage{tabularx}
\usepackage{enumitem}
\usepackage{xcolor}
\usepackage{csquotes}
\usepackage{float}
\usepackage{tikz}
\usetikzlibrary{arrows.meta,positioning,shapes.geometric,calc,fit,backgrounds}
\usepackage{pgfplots}
\pgfplotsset{compat=1.18}
\usepackage{xurl}
\usepackage[hidelinks]{hyperref}
\usepackage[backend=biber,style=apa,sorting=nyt]{biblatex}
\DeclareLanguageMapping{spanish}{spanish-apa}

\geometry{margin=2.5cm}
\setlength{\emergencystretch}{1em}
\setlength{\parindent}{0pt}
\setlength{\parskip}{0.65em}
\setlist{nosep}
\addbibresource{referencias.bib}
\AtBeginBibliography{\sloppy}

% --------------------------------------------------------------------------
% Estilos gráficos comunes
% --------------------------------------------------------------------------
\definecolor{capa}{HTML}{2F5D8C}
\definecolor{trans}{HTML}{8C5A2F}
\definecolor{acento}{HTML}{9C2B2B}

\tikzset{
  % Notación OPM
  opmobj/.style   = {draw, very thick, rectangle, fill=green!8,
                     font=\footnotesize\bfseries, align=center,
                     inner sep=4pt, minimum height=8mm, text width=26mm},
  opmobjn/.style  = {opmobj, text width=}, % objeto de ancho natural
  opmproc/.style  = {draw, very thick, ellipse, fill=blue!12,
                     font=\footnotesize\bfseries, align=center,
                     inner sep=1pt, minimum height=9mm, minimum width=28mm},
  opmstate/.style = {draw, rounded corners=4pt, fill=white, font=\scriptsize,
                     inner sep=2pt, minimum height=4.5mm},
  instr/.style    = {thick, -{Circle[open,length=5pt]}},
  agente/.style   = {thick, -{Circle[length=5pt]}},
  efecto/.style   = {thick, {Stealth[length=5.5pt]}-{Stealth[length=5.5pt]}},
  flecha/.style   = {thick, -{Stealth[length=5pt]}},
  % Diagramas de bloques
  bloque/.style   = {draw, thick, rounded corners=2pt, align=center,
                     font=\footnotesize, inner sep=5pt},
  etiq/.style     = {font=\scriptsize\itshape, align=center},
}

% Triángulos estructurales de OPM.
% #1 = rotación en grados (0 = base arriba, vértice abajo)
% #2 = nombre de una coordenada previamente definida
\newcommand{\opmagg}[2]{% agregación: triángulo relleno
  \begin{scope}[rotate around={#1:(#2)}]
    \path[draw,very thick,fill=black] ([shift={(-2mm,2mm)}]#2) --
      ([shift={(2mm,2mm)}]#2) -- ([shift={(0mm,-2mm)}]#2) -- cycle;
  \end{scope}}
\newcommand{\opmgen}[2]{% generalización: triángulo hueco
  \begin{scope}[rotate around={#1:(#2)}]
    \path[draw,very thick,fill=white] ([shift={(-2mm,2mm)}]#2) --
      ([shift={(2mm,2mm)}]#2) -- ([shift={(0mm,-2mm)}]#2) -- cycle;
  \end{scope}}
\newcommand{\opmexh}[2]{% exhibición: triángulo hueco con triángulo relleno
  \begin{scope}[rotate around={#1:(#2)}]
    \path[draw,very thick,fill=white] ([shift={(-2.4mm,2.4mm)}]#2) --
      ([shift={(2.4mm,2.4mm)}]#2) -- ([shift={(0mm,-2.4mm)}]#2) -- cycle;
    \path[fill=black] ([shift={(-1mm,1.5mm)}]#2) --
      ([shift={(1mm,1.5mm)}]#2) -- ([shift={(0mm,-0.1mm)}]#2) -- cycle;
  \end{scope}}

\title{\textbf{Alcance de la tecnología: la infraestructura como código}\\
\large Funcionamiento de un sistema completo, capas y módulos principales,
y modelo de objetos y procesos}
\author{Carlos Luis Rojas Aragonés\\
Gestión del Desarrollo Tecnológico}
\date{27 de agosto de 2026}

\begin{document}

\maketitle

\begin{abstract}
\noindent\small
La infraestructura como código (\emph{Infrastructure as Code}, IaC) sustituye
la configuración manual de servidores, redes y servicios por una
especificación versionada que un motor automático traduce en operaciones sobre
la interfaz de programación de un proveedor. Este trabajo delimita el alcance
de la tecnología, describe su principio de funcionamiento (la conciliación
continua entre un estado deseado y un estado real), desglosa las siete capas y
los seis módulos transversales de un sistema completo, y formaliza el conjunto
mediante un modelo de objetos y procesos según la metodología OPM, con su
diagrama de sistema, sus enlaces estructurales, su ampliación en detalle y el
lenguaje OPL correspondiente. Todas las figuras son de elaboración propia.
\end{abstract}

% ==========================================================================
\section{Objeto y delimitación del alcance}
% ==========================================================================

\subsection*{01. El encargo}

Esta primera parte de la actividad pide tres cosas: un resumen del
funcionamiento básico de un sistema completo de infraestructura como código,
un conjunto de imágenes y diagramas que expliquen sus capas o módulos
principales, y un diagrama de proceso-objeto elaborado en el entorno OPM
Sandbox \parencite{opcloud}. El documento atiende las dos primeras en las
secciones 2 a 8, con el resumen funcional y los diagramas de capas y módulos,
y la tercera en la sección 9, que contiene el modelo de objetos y procesos con
su diagrama de sistema, sus enlaces estructurales, su ampliación en detalle y
el lenguaje OPL correspondiente. El diagrama construido en el propio entorno
de la herramienta, que no guarda ni exporta el trabajo, se entrega aparte.

\subsection*{02. Qué se entiende aquí por sistema completo}

Conviene fijar el término antes de dibujar nada, porque \enquote{infraestructura
como código} se usa con dos alcances muy distintos. En el sentido estrecho
designa una herramienta: un fichero declarativo y un binario que lo aplica. En
el sentido amplio, que es el que aquí interesa, designa un sistema
sociotécnico completo, con su cadena de herramientas, su modelo de estado, sus
controles de gobierno y el equipo que lo opera.

Un sistema completo de IaC es aquel que satisface cuatro propiedades
simultáneas \parencite{morris2021}:

\begin{enumerate}
  \item \textbf{Especificación externalizada.} La configuración deseada de la
    infraestructura vive fuera de la infraestructura, en artefactos de texto
    legibles y versionados, no en la memoria de un operador ni en el estado
    acumulado de las máquinas.
  \item \textbf{Reproducibilidad.} Aplicar la misma especificación sobre un
    sustrato vacío produce un resultado equivalente, con independencia de
    quién la aplique y de cuándo lo haga.
  \item \textbf{Idempotencia.} Aplicar la especificación sobre un sistema que
    ya la cumple no produce ningún cambio, y aplicarla \emph{n} veces tiene el
    mismo efecto que aplicarla una vez.
  \item \textbf{Conciliación.} El sistema es capaz de detectar la divergencia
    entre lo especificado y lo realmente desplegado, y de corregirla.
\end{enumerate}

La literatura académica sobre la materia, todavía escasa en comparación con su
difusión industrial, se ha concentrado sobre todo en la calidad y los defectos
del propio código de infraestructura \parencite{rahman2019}.

Un montaje que cumpla sólo la primera propiedad es un repositorio de guiones,
no un sistema de IaC. La cuarta es la que convierte la automatización en
control, y es también la que con más frecuencia falta en las implantaciones
reales.

\subsection*{03. La frontera del sistema}

Delimitar el alcance exige decir también qué queda fuera. El
Cuadro~\ref{tab:frontera} fija la frontera adoptada en este documento.

\begin{table}[ht]
  \centering
  \footnotesize
  \caption{Frontera del sistema de infraestructura como código.}
  \label{tab:frontera}
  \begin{tabularx}{\textwidth}{@{}>{\raggedright\arraybackslash}X
    >{\raggedright\arraybackslash}X@{}}
    \toprule
    \textbf{Dentro del alcance} & \textbf{Fuera del alcance} \\
    \midrule
    Cómputo, red, almacenamiento, identidad y servicios gestionados
      declarados como código. &
    El sustrato físico del proveedor: centros de datos, cableado,
      refrigeración. \\
    Configuración del sistema operativo y de los agentes instalados en él. &
    El código de la aplicación de negocio y sus pruebas funcionales. \\
    Construcción de imágenes y artefactos inmutables. &
    Los datos de producción y su ciclo de vida, que exigen migraciones y
      copias de seguridad con reglas propias. \\
    Orquestación de cargas y despliegue continuo de manifiestos. &
    La gestión de incidencias, la guardia y los procedimientos de crisis. \\
    Estado, secretos, políticas, pruebas y observabilidad de la propia
      infraestructura. &
    La observabilidad funcional de la aplicación (trazas de negocio,
      experiencia de usuario). \\
    \bottomrule
  \end{tabularx}
\end{table}

La distinción entre infraestructura y datos merece un comentario, porque es la
frontera que más problemas causa en la práctica. Un servidor de base de datos
está dentro del alcance: su tamaño, su versión, su red y sus copias
programadas se declaran como código. El contenido de esa base de datos no lo
está: no es idempotente, no se puede destruir y recrear, y su presencia obliga
a que el sistema de IaC distinga recursos desechables de recursos con estado
persistente, y proteja a los segundos de las operaciones de reemplazo.

% ==========================================================================
\section{El principio de funcionamiento}
% ==========================================================================

\subsection*{01. Tres estados y una función}

Todo el funcionamiento de un sistema de IaC se sostiene sobre tres
representaciones de la misma infraestructura y sobre la función que las
compara.

\begin{description}[style=nextline, leftmargin=0pt, labelwidth=0pt, labelsep=0pt]
  \item[Estado deseado.] Lo que la especificación declara. Es un artefacto de
    texto, versionado, que describe recursos y sus propiedades. No describe
    los pasos para llegar a ellos.
  \item[Estado registrado.] Lo que el motor cree que existe, resultado de la
    última operación que él mismo ejecutó. Se materializa en un fichero o
    servicio de estado que asocia cada recurso declarado con el identificador
    real que el proveedor le asignó \parencite{terraformstate}.
  \item[Estado real.] Lo que efectivamente existe en el proveedor, que sólo se
    conoce consultando su interfaz de programación.
\end{description}

La función central del sistema es la \textbf{conciliación}: calcular la
diferencia entre el estado deseado y el estado real, y emitir el conjunto
mínimo de operaciones que anula esa diferencia. La Figura~\ref{fig:bucle}
representa el bucle.

\begin{figure}[ht]
  \centering
  \begin{tikzpicture}
    \node[bloque, fill=blue!8, text width=30mm, minimum height=12mm] (des)
      {\textbf{Estado deseado}\\\scriptsize especificación versionada};
    \node[bloque, fill=orange!12, text width=30mm, minimum height=12mm,
      right=26mm of des] (comp)
      {\textbf{Comparador}\\\scriptsize cálculo de la diferencia};
    \node[bloque, fill=green!10, text width=30mm, minimum height=12mm,
      right=26mm of comp] (real)
      {\textbf{Estado real}\\\scriptsize recursos del proveedor};
    \node[bloque, fill=gray!10, text width=30mm, minimum height=12mm,
      below=17mm of comp] (act)
      {\textbf{Actuador}\\\scriptsize crear, modificar, destruir};

    \draw[flecha] (des) -- node[above, etiq] {declaración} (comp);
    \draw[flecha] (real) -- node[above, etiq] {observación} (comp);
    \draw[flecha] (comp) -- node[right, etiq, xshift=1mm] {plan} (act);
    \draw[flecha] (act.east) to[out=0, in=-90]
      node[pos=0.55, below right, etiq, xshift=-3mm, yshift=1mm]
      {\parbox{24mm}{\centering operaciones\\sobre la API}} (real.south);
  \end{tikzpicture}
  \caption{Bucle de conciliación. El sistema no ejecuta una secuencia de
  pasos, sino que cierra repetidamente la distancia entre dos
  representaciones. Si la diferencia es vacía, el actuador no ejecuta nada, y
  ésa es exactamente la propiedad de idempotencia. Elaboración propia.}
  \label{fig:bucle}
\end{figure}

\subsection*{02. Declarativo frente a imperativo}

La diferencia entre los dos paradigmas no es de sintaxis sino de quién calcula
el camino. En el enfoque imperativo, el autor escribe la secuencia de
operaciones y asume la responsabilidad de que sea correcta desde cualquier
estado de partida. En el enfoque declarativo, el autor escribe el destino y el
motor deriva la secuencia.

La consecuencia práctica es que un guion imperativo sólo es correcto para el
estado de partida que su autor imaginó, mientras que una especificación
declarativa es correcta para cualquiera, a condición de que el motor sepa
observar el estado real. Por eso el paradigma declarativo domina la capa de
aprovisionamiento, donde el espacio de estados de partida es enorme, y
sobrevive el imperativo en tareas de configuración muy locales, donde ese
espacio es pequeño y conocido.

Casi ningún sistema es puro. Las herramientas de configuración expresan los
pasos en orden, pero cada paso es un módulo idempotente que sólo actúa si hace
falta \parencite{ansibledocs}. Es un híbrido: secuencia imperativa de
operaciones declarativas.

\subsection*{03. Convergencia y deriva}

Un sistema convergente es el que, aplicado repetidamente, se aproxima al
estado deseado y luego se queda quieto. La \textbf{deriva}
(\emph{drift}) es la aparición de diferencias entre el estado real y el
deseado por causas ajenas al sistema: un cambio manual de urgencia, una acción
automática del proveedor, la caducidad de un certificado, el borrado
accidental de un recurso.

La deriva es inevitable, y la madurez de una implantación se mide por lo que
hace con ella. Hay tres posturas, en orden creciente de rigor:

\begin{enumerate}
  \item \textbf{Ignorarla.} El motor sólo se ejecuta cuando alguien lo lanza.
    Entre ejecuciones, el estado real puede alejarse arbitrariamente.
  \item \textbf{Detectarla.} Una ejecución periódica en modo de sólo lectura
    compara y notifica, sin corregir. Convierte la deriva en una métrica
    observable.
  \item \textbf{Corregirla.} Un agente concilia de forma continua y revierte
    todo cambio no declarado. Es el principio de conciliación continua de
    GitOps \parencite{opengitops} y el modo nativo de los controladores de
    Kubernetes \parencite{k8scontrollers,burns2016}.
\end{enumerate}

% ==========================================================================
\section{El ciclo básico de operación}
% ==========================================================================

Descendiendo un nivel desde el bucle abstracto, la operación cotidiana de un
sistema completo recorre siete pasos. La Figura~\ref{fig:ciclo} los ordena.

\begin{figure}[ht]
  \centering
  \begin{tikzpicture}[node distance=4.5mm and 5.5mm]
    \tikzset{paso/.style={bloque, fill=capa!12, draw=capa, text width=20mm,
      minimum height=13mm, font=\scriptsize}}
    \node[paso] (p1) {\textbf{1. Codificar}\\ escribir la\\especificación};
    \node[paso, right=of p1] (p2) {\textbf{2. Validar}\\ formato, tipos,\\
      políticas};
    \node[paso, right=of p2] (p3) {\textbf{3. Planificar}\\ diferencia\\
      deseado--real};
    \node[paso, right=of p3] (p4) {\textbf{4. Autorizar}\\ revisión y\\
      aprobación};
    \node[paso, right=of p4] (p5) {\textbf{5. Aplicar}\\ ejecutar sobre\\
      la API};
    \node[paso, below=9mm of p4] (p6) {\textbf{6. Verificar}\\ pruebas de\\
      infraestructura};
    \node[paso, below=9mm of p2] (p7) {\textbf{7. Conciliar}\\ vigilar y\\
      corregir la deriva};

    \foreach \a/\b in {p1/p2,p2/p3,p3/p4,p4/p5}
      \draw[flecha, draw=capa] (\a) -- (\b);
    \draw[flecha, draw=capa] (p5.south) |- (p6.east);
    \draw[flecha, draw=capa] (p6) -- (p7);
    \draw[flecha, draw=acento, dashed] (p7.west) -|
      node[pos=0.28, below, etiq, text=acento, yshift=-1mm]
      {deriva detectada} (p1.south);
  \end{tikzpicture}
  \caption{Los siete pasos del ciclo de operación. El retorno discontinuo es
  el que distingue un sistema de IaC de una simple automatización de
  despliegue. Elaboración propia.}
  \label{fig:ciclo}
\end{figure}

Los pasos 2, 4 y 6 son los que suelen faltar en implantaciones incompletas, y
son precisamente los que aportan garantía. Sin validación de políticas, el
sistema automatiza también los errores de configuración; sin autorización,
elimina el punto de control humano sobre cambios irreversibles; sin
verificación, no distingue entre \enquote{el motor terminó sin error} y
\enquote{la infraestructura funciona}.

\subsection*{01. El paso crítico: la comparación a tres bandas}

El paso 3 merece detalle porque en él se concentra el mecanismo que hace útil
al sistema. El motor no compara dos cosas sino tres, y lo hace en dos fases
(Figura~\ref{fig:tresbandas}).

\begin{figure}[ht]
  \centering
  \begin{tikzpicture}
    \tikzset{caja/.style={bloque, text width=34mm, minimum height=15mm}}
    \node[caja, fill=blue!8]  (d) {\textbf{Deseado}\\\scriptsize lo que
      declara el código};
    \node[caja, fill=gray!12, right=20mm of d] (r) {\textbf{Registrado}\\
      \scriptsize lo que el motor anotó\\ en la última ejecución};
    \node[caja, fill=green!10, right=20mm of r] (a) {\textbf{Real}\\
      \scriptsize lo que devuelve\\ la API del proveedor};

    \draw[efecto, draw=trans] (r) -- node[above, etiq, text=trans]
      {\textbf{fase 1}} node[below, etiq, text=trans] {refresco} (a);
    \draw[efecto, draw=acento] (d) -- node[above, etiq, text=acento]
      {\textbf{fase 2}} node[below, etiq, text=acento] {plan} (r);

    \node[bloque, fill=acento!10, draw=acento, below=12mm of r,
      text width=0.62\textwidth]
      {\textbf{Plan de cambios}: lista de recursos a crear, modificar,
       reemplazar o destruir, con el valor anterior y el posterior de cada
       propiedad afectada};
    \draw[flecha, draw=acento] ($(d.south)!0.5!(a.south)+(0,-2mm)$) --
      ++(0,-5mm);
  \end{tikzpicture}
  \caption{Comparación a tres bandas. La fase de refresco actualiza el
  registro contra el mundo real y es la que descubre la deriva; la fase de
  plan compara el código contra ese registro ya actualizado y es la que
  produce las operaciones. Elaboración propia.}
  \label{fig:tresbandas}
\end{figure}

El estado registrado no es redundante. Cumple tres funciones que ninguna de
las otras dos representaciones puede cumplir. Primera, traduce el nombre
lógico que el código da a un recurso al identificador opaco que el proveedor
le asignó, sin lo cual el motor no podría reconocer que un recurso existente
le pertenece. Segunda, permite detectar destrucciones: un recurso presente en
el registro y ausente del proveedor sólo puede detectarse comparando contra el
registro, nunca comparando el código con la realidad. Tercera, guarda
propiedades calculadas por el proveedor que el código no declara y que otros
recursos necesitan.

Ese fichero es, en consecuencia, un activo crítico: contiene identificadores y
con frecuencia valores sensibles, y su corrupción o pérdida deja al sistema
sin capacidad de reconocer lo que él mismo creó. De ahí que todo montaje serio
lo almacene en un servicio remoto con cifrado, versionado y bloqueo de
concurrencia.

% ==========================================================================
\section{Las capas de un sistema completo}
% ==========================================================================

Un sistema completo se organiza en siete capas apiladas, cada una de las
cuales consume el resultado de la anterior y expone una abstracción más alta.
La Figura~\ref{fig:pila} las presenta con sus herramientas representativas.

\begin{figure}[ht]
  \centering
  \begin{tikzpicture}[y=13.6mm]
    \def\W{14.6}
    \foreach \i/\nom/\fun/\tool in {
      0/{L0 · Sustrato}/{cómputo, red, almacenamiento e identidad del
         proveedor}/{API de AWS, Azure o GCP; equipo físico},
      1/{L1 · Aprovisionamiento}/{crea, modifica y destruye los recursos
         declarados}/{Terraform, OpenTofu, Pulumi, CloudFormation, Bicep},
      2/{L2 · Imágenes}/{construye artefactos inmutables ya
         configurados}/{Packer, Dockerfile, \emph{buildpacks}},
      3/{L3 · Configuración}/{ajusta el interior de la máquina ya
         arrancada}/{Ansible, Chef, Puppet, cloud-init},
      4/{L4 · Orquestación}/{programa cargas sobre el conjunto de
         máquinas}/{Kubernetes, Nomad, Helm, Kustomize},
      5/{L5 · Entrega}/{lleva el código al motor y concilia de forma
         continua}/{Argo CD, Flux, canalizaciones de integración},
      6/{L6 · Plataforma}/{expone catálogo y autoservicio al equipo de
         producto}/{Backstage, Crossplane, portales internos}} {
      \node[draw=capa, very thick, fill=capa!\the\numexpr 6+2*\i\relax,
        rounded corners=2pt, minimum width=\W cm, minimum height=12.6mm,
        anchor=south west] (L\i) at (0,\i) {};
      \node[anchor=west, font=\footnotesize\bfseries, text=capa!60!black,
        text width=37mm, align=left] at ([xshift=3.5mm]L\i.west) {\nom};
      \node[anchor=west, font=\scriptsize, text width=50mm, align=left]
        at ([xshift=43mm]L\i.west) {\fun};
      \node[anchor=west, font=\scriptsize\itshape, text width=45mm,
        align=left] at ([xshift=96mm]L\i.west) {\tool};
    }
    \draw[flecha, draw=capa, line width=1pt] (-0.5,0.2) -- (-0.5,6.8)
      node[midway, rotate=90, above, font=\scriptsize, text=capa]
      {abstracción creciente};
  \end{tikzpicture}
  \caption{Las siete capas de un sistema completo de infraestructura como
  código. Cada capa consume la salida de la inferior. Las herramientas citadas
  son representativas, no exhaustivas. Elaboración propia.}
  \label{fig:pila}
\end{figure}

\subsection*{01. Qué hace cada capa}

\textbf{L0, sustrato.} No es parte del sistema de IaC sino su objeto. Aporta
la interfaz de programación sobre la que todo lo demás actúa. Su
característica determinante es que ofrece recursos con identidad propia,
creables y destruibles por llamada, lo que es la condición de posibilidad de
todo el resto.

\textbf{L1, aprovisionamiento.} Traduce la especificación en llamadas de
creación, modificación y destrucción. Es la capa que posee el registro de
estado y la que ejecuta la comparación a tres bandas. Trabaja con recursos que
todavía no tienen sistema operativo arrancado, o cuyo interior no le
concierne. Conviene señalar que en 2023 el cambio de licencia de Terraform dio
origen a OpenTofu, una bifurcación acogida por la Linux Foundation
\parencite{opentofu}, de modo que hoy la capa cuenta con dos implementaciones
compatibles del mismo lenguaje: una decisión de dependencia que una hoja de
ruta tecnológica debería registrar de forma explícita.

\textbf{L2, imágenes.} Desplaza el trabajo de configuración desde el momento
de ejecución al momento de construcción. En lugar de arrancar una máquina
genérica y modificarla, construye una imagen ya completa y arranca instancias
idénticas de ella. Es la base de la infraestructura inmutable: en lugar de
modificar un servidor, se reemplaza. Reduce drásticamente la deriva, porque
elimina la ventana temporal en la que un servidor puede divergir.

\textbf{L3, configuración.} Actúa sobre el interior de máquinas ya arrancadas:
paquetes, ficheros, servicios, usuarios. Es la capa históricamente más antigua
y la que más terreno ha cedido a L2 y L4. Sigue siendo necesaria donde hay
máquinas de larga vida, sistemas heredados o equipos físicos.

\textbf{L4, orquestación.} Introduce una indirección importante: deja de
declararse \enquote{esta carga en esta máquina} y pasa a declararse
\enquote{esta carga, tantas réplicas} sobre un conjunto de máquinas
intercambiables. El planificador decide la colocación. Sus controladores son
bucles de conciliación como el de la Figura~\ref{fig:bucle}, pero embebidos y
permanentes \parencite{k8scontrollers}.

\textbf{L5, entrega.} Conecta el repositorio con el motor. Ejecuta las
canalizaciones que validan, planifican y aplican, y, en el modelo de
conciliación continua, aloja el agente que vigila la deriva sin intervención
humana \parencite{opengitops,humble2010}.

\textbf{L6, plataforma.} Envuelve todo lo anterior en una interfaz de
autoservicio, para que un equipo de producto pueda solicitar una base de datos
sin conocer ni la capa L1 ni las políticas de la organización, que quedan
codificadas en el catálogo.

\subsection*{02. La tensión entre capas}

Las capas L2, L3 y L4 compiten parcialmente por el mismo territorio. Un mismo
requisito, por ejemplo \enquote{que este servidor tenga la versión 1.24 de un
tiempo de ejecución}, puede resolverse construyendo una imagen que ya la
incluya (L2), instalándola con un módulo de configuración (L3) o empaquetando
la aplicación en un contenedor que la lleve dentro (L4). Las tres soluciones
funcionan y no deben coexistir sobre el mismo recurso: si dos capas gobiernan
la misma propiedad, cada una revierte el trabajo de la otra en cada ciclo de
conciliación, y el sistema oscila.

Decidir qué capa posee cada propiedad es, por tanto, la primera decisión de
arquitectura de una implantación, y conviene documentarla de forma explícita.

% ==========================================================================
\section{Los módulos transversales}
% ==========================================================================

Las siete capas describen el flujo principal, pero un sistema completo
necesita además seis módulos que atraviesan todas las capas
(Figura~\ref{fig:transversales}). Ninguno produce infraestructura; todos
condicionan si la que se produce es correcta, segura y auditable.

\begin{figure}[ht]
  \centering
  \begin{tikzpicture}
    \node[draw=capa, very thick, fill=capa!10, rounded corners=3pt,
      minimum width=5.4cm, minimum height=3.7cm, align=center,
      font=\footnotesize\bfseries, text=capa!60!black] (nucleo)
      {Las siete capas\\del flujo principal\\[2pt]
       \scriptsize\mdseries L0 sustrato · L1 aprovisionamiento\\
       \scriptsize\mdseries L2 imágenes · L3 configuración\\
       \scriptsize\mdseries L4 orquestación · L5 entrega\\
       \scriptsize\mdseries L6 plataforma};
    \foreach \ang/\nom/\det in {
      90/{Control de versiones}/{historia, revisión, reversión},
      32/{Gestión de estado}/{registro, bloqueo, cifrado},
      -32/{Gestión de secretos}/{inyección en tiempo de ejecución},
      -90/{Política como código}/{reglas evaluables antes de aplicar},
      -148/{Pruebas}/{estáticas, unitarias, de integración},
      148/{Observabilidad}/{deriva, coste, cumplimiento}} {
      \node[bloque, fill=trans!12, draw=trans, text width=31mm,
        font=\scriptsize, align=center, anchor=center] at (\ang:5.1cm)
        (m\ang) {\textbf{\nom}\\ \det};
      \draw[efecto, draw=trans, shorten >=3pt, shorten <=3pt]
        (m\ang) -- (nucleo);
    }
  \end{tikzpicture}
  \caption{Los seis módulos transversales. Cada uno interactúa con todas las
  capas del flujo principal, no con una en particular. Elaboración propia.}
  \label{fig:transversales}
\end{figure}

El Cuadro~\ref{tab:transversales} resume la función de cada uno y el fallo
característico que aparece cuando falta.

\begin{table}[ht]
  \centering
  \footnotesize
  \caption{Módulos transversales: función, instrumentos y fallo característico
  por ausencia.}
  \label{tab:transversales}
  \begin{tabularx}{\textwidth}{@{}>{\raggedright\arraybackslash}p{0.145\textwidth}
    >{\raggedright\arraybackslash}X
    >{\raggedright\arraybackslash}p{0.30\textwidth}@{}}
    \toprule
    \textbf{Módulo} & \textbf{Función} & \textbf{Fallo por ausencia} \\
    \midrule
    Control de versiones & Registrar quién cambió qué, cuándo y por qué;
      permitir revertir a un estado anterior conocido. &
      Cambios sin trazabilidad y reversión imposible. \\
    Gestión de estado & Custodiar el registro: almacenamiento remoto, cifrado,
      versionado y bloqueo para impedir ejecuciones concurrentes. &
      Corrupción del registro y recursos huérfanos. \\
    Gestión de secretos & Mantener credenciales fuera del código e inyectarlas
      en el momento de la ejecución, con rotación y ámbito mínimo. Es uno de
      los defectos más frecuentes en código de infraestructura
      \parencite{rahman2019smells}. &
      Credenciales escritas en el repositorio. \\
    Política como código & Expresar las reglas de la organización como
      predicados evaluables sobre el plan, antes de aplicarlo
      \parencite{opa}. &
      El sistema automatiza también los incumplimientos, y a mayor
      velocidad. \\
    Pruebas & Verificar la especificación en varios niveles: análisis
      estático, validación del plan, despliegue efímero y comprobación
      funcional. &
      \enquote{Terminó sin error} se confunde con \enquote{funciona}. \\
    Observabilidad & Medir deriva, coste, tiempo de ciclo y cumplimiento, y
      exponerlos como series temporales. &
      No hay figuras de mérito y la mejora no es demostrable. \\
    \bottomrule
  \end{tabularx}
\end{table}

% ==========================================================================
\section{Modos de operación}
% ==========================================================================

Sobre la misma arquitectura caben decisiones de operación que cambian de forma
apreciable el comportamiento del sistema. Las tres principales son
independientes entre sí.

\subsection*{01. Empuje frente a arrastre}

En el modo de \textbf{empuje} (\emph{push}), un actor externo, típicamente una
canalización de integración, posee las credenciales del entorno y ejecuta el
motor contra él. En el modo de \textbf{arrastre} (\emph{pull}), un agente
residente dentro del entorno consulta el repositorio y se aplica a sí mismo
los cambios. La Figura~\ref{fig:pushpull} contrasta ambos.

\begin{figure}[ht]
  \centering
  \begin{tikzpicture}
    \tikzset{c/.style={bloque, text width=28mm, minimum height=9mm},
             cin/.style={bloque, fill=white, text width=30mm,
               minimum height=7mm, font=\scriptsize},
             ent/.style={draw, thick, rounded corners=2pt, fill=green!10,
               minimum width=52mm}}
    % ---------------- empuje ----------------
    \node[etiq, font=\footnotesize\bfseries] (t1) {Empuje (\emph{push})};
    \node[c, fill=blue!8, below=5mm of t1] (r1) {Repositorio};
    \node[c, fill=orange!12, below=11mm of r1] (ci1)
      {Canalización\\ \scriptsize con las credenciales};
    \node[ent, minimum height=24mm, below=11mm of ci1] (e1) {};
    \node[anchor=north, font=\scriptsize\bfseries]
      at ([yshift=-2mm]e1.north) {Entorno};
    \node[cin, anchor=south] at ([yshift=3mm]e1.south) (i1)
      {Infraestructura};
    \draw[flecha] (r1) -- node[right, etiq, xshift=1mm] {disparo} (ci1);
    \draw[flecha] (ci1) -- node[right, etiq, xshift=1mm] {aplica} (e1.north);
    % ---------------- arrastre ----------------
    \node[etiq, font=\footnotesize\bfseries, right=52mm of t1] (t2)
      {Arrastre (\emph{pull})};
    \node[c, fill=blue!8, below=5mm of t2] (r2) {Repositorio};
    \node[ent, minimum height=34mm, below=26mm of r2] (e2) {};
    \node[anchor=north, font=\scriptsize\bfseries]
      at ([yshift=-1.5mm]e2.north) {Entorno};
    \node[cin, fill=orange!12, anchor=north] at ([yshift=-7mm]e2.north) (ag)
      {Agente\\ \scriptsize con las credenciales};
    \node[cin, anchor=south] at ([yshift=3mm]e2.south) (i2)
      {Infraestructura};
    \draw[flecha] ([xshift=-11mm]e2.north) to[out=95,in=-95]
      node[left, etiq, xshift=-1mm] {consulta} ([xshift=-9mm]r2.south);
    \draw[flecha] ([xshift=9mm]r2.south) to[out=-85,in=85]
      node[right, etiq, xshift=1mm] {\parbox{18mm}{\centering estado\\
      deseado}} ([xshift=11mm]e2.north);
    \draw[flecha, draw=acento] (ag) -- node[right, font=\tiny\itshape,
      text=acento, xshift=0.5mm] {concilia} (i2);
  \end{tikzpicture}
  \caption{Empuje frente a arrastre. En el arrastre, las credenciales del
  entorno nunca salen de él, lo que reduce la superficie de ataque y hace del
  agente el punto natural de conciliación continua. Elaboración propia a
  partir de los principios de \textcite{opengitops}.}
  \label{fig:pushpull}
\end{figure}

El arrastre tiene dos ventajas estructurales. La primera es de seguridad: las
credenciales de producción no viajan a un sistema de integración compartido.
La segunda es que el agente, por estar siempre presente, puede conciliar de
forma continua y no sólo cuando alguien confirma un cambio, que es la cuarta
propiedad exigida en la sección 1. Su coste es que exige un agente por
entorno y complica los flujos que necesitan orden entre entornos.

\subsection*{02. Mutable frente a inmutable}

En el modelo mutable, un recurso desplegado se modifica en su sitio. En el
inmutable, no se modifica nunca: se construye uno nuevo con la configuración
deseada y se sustituye. El modelo inmutable elimina por construcción la deriva
de configuración y hace que el recurso desplegado sea siempre trazable a un
artefacto concreto. A cambio, obliga a que todo estado persistente viva fuera
del recurso reemplazable, lo que es una restricción de diseño exigente.

\subsection*{03. El cuadro de decisión}

Las tres decisiones son independientes entre sí, de modo que caben ocho
combinaciones. La más frecuente en implantaciones recientes es la de arrastre,
inmutable y declarativa, pero ninguna de las tres opciones es obligada. El
Cuadro~\ref{tab:modos} resume las consecuencias de cada una.

\begin{table}[!ht]
  \centering
  \footnotesize
  \caption{Tres decisiones de operación y sus consecuencias.}
  \label{tab:modos}
  \begin{tabularx}{\textwidth}{@{}>{\raggedright\arraybackslash}p{0.15\textwidth}
    >{\raggedright\arraybackslash}X>{\raggedright\arraybackslash}X@{}}
    \toprule
    \textbf{Decisión} & \textbf{Opción A} & \textbf{Opción B} \\
    \midrule
    Iniciativa &
      \textbf{Empuje}: control de orden entre entornos; credenciales fuera del
      entorno; conciliación sólo bajo disparo. &
      \textbf{Arrastre}: credenciales confinadas; conciliación continua;
      un agente por entorno. \\
    Mutabilidad &
      \textbf{Mutable}: cambios rápidos y baratos; deriva acumulativa; difícil
      trazar qué versión corre. &
      \textbf{Inmutable}: sin deriva; reemplazo trazable; obliga a externalizar
      todo estado persistente. \\
    Expresión &
      \textbf{Declarativa}: el motor calcula el camino; correcta desde
      cualquier estado de partida. &
      \textbf{Imperativa}: control fino del orden; sólo correcta desde el
      estado de partida previsto. \\
    \bottomrule
  \end{tabularx}
\end{table}

% ==========================================================================
\section{El módulo como unidad de composición}
% ==========================================================================

El término \enquote{módulo} tiene en IaC un sentido preciso: es la unidad
reutilizable que encapsula un conjunto de recursos tras una interfaz de
entradas y salidas. Es el mecanismo que permite que una organización no repita
la misma configuración en cada equipo y, sobre todo, que las decisiones
correctas queden incorporadas por omisión.

\begin{figure}[ht]
  \centering
  \begin{tikzpicture}
    \node[draw=capa, very thick, fill=capa!8, rounded corners=3pt,
      minimum width=5.6cm, minimum height=3.2cm] (mod) {};
    \node[anchor=north, font=\footnotesize\bfseries, text=capa!60!black]
      at ([yshift=-3mm]mod.north) {Módulo \texttt{red-privada v2.3.1}};
    \node[bloque, fill=white, font=\scriptsize, text width=44mm,
      anchor=north] at ([yshift=-9mm]mod.north)
      {\textbf{Cuerpo}: recursos declarados,\\ valores por omisión,
       validaciones,\\ etiquetas obligatorias};

    \node[bloque, fill=blue!8, text width=27mm, font=\scriptsize,
      left=16mm of mod] (in)
      {\textbf{Entradas}\\ región, CIDR,\\ número de zonas,\\ etiquetas};
    \node[bloque, fill=green!10, text width=27mm, font=\scriptsize,
      right=16mm of mod] (out)
      {\textbf{Salidas}\\ identificador de red,\\ subredes,\\ tabla de rutas};
    \node[bloque, fill=trans!12, text width=40mm, font=\scriptsize,
      below=9mm of mod] (reg)
      {\textbf{Registro versionado}\\ publica, firma y resuelve versiones};

    \draw[flecha] (in) -- (mod);
    \draw[flecha] (mod) -- (out);
    \draw[flecha] (reg) -- (mod);

    \node[bloque, fill=gray!8, text width=25mm, font=\scriptsize,
      above left=8mm and 8mm of mod.north] (c1) {Composición\\
      \emph{entorno de pruebas}};
    \node[bloque, fill=gray!8, text width=25mm, font=\scriptsize,
      above right=8mm and 8mm of mod.north] (c2) {Composición\\
      \emph{entorno de producción}};
    \draw[flecha, gray] (c1) -- (mod.north west);
    \draw[flecha, gray] (c2) -- (mod.north east);
  \end{tikzpicture}
  \caption{Anatomía de un módulo. Las mismas entradas con valores distintos
  producen entornos comparables; el registro versionado permite que una
  corrección se propague de forma controlada. Elaboración propia.}
  \label{fig:modulo}
\end{figure}

Tres propiedades hacen que un módulo cumpla su función. Debe tener una
\textbf{interfaz estrecha}: cuantas menos entradas expone, más decisiones fija
y más valor aporta. Debe estar \textbf{versionado} con semántica explícita,
para que quien lo consume elija cuándo adoptar un cambio. Y debe delimitar un
\textbf{radio de explosión} razonable: un módulo que agrupa demasiados
recursos obliga a que cada cambio pequeño planifique sobre un conjunto grande,
lo que alarga el ciclo y aumenta el riesgo de cada aplicación.

% ==========================================================================
\section{Figuras de mérito del sistema}
% ==========================================================================

Un sistema de IaC es en sí mismo una tecnología, y como tal admite figuras de
mérito que permiten situarlo y seguir su evolución. El
Cuadro~\ref{tab:fom} propone ocho, con su unidad y su procedimiento de
medida. Las cuatro primeras coinciden con las métricas de rendimiento de
entrega ampliamente utilizadas \parencite{forsgren2018}; las cuatro últimas
son específicas de la infraestructura.

\begin{table}[ht]
  \centering
  \footnotesize
  \caption{Figuras de mérito propuestas para un sistema de IaC.}
  \label{tab:fom}
  \begin{tabularx}{\textwidth}{@{}>{\raggedright\arraybackslash}p{0.27\textwidth}
    >{\raggedright\arraybackslash}p{0.17\textwidth}
    >{\raggedright\arraybackslash}X@{}}
    \toprule
    \textbf{Figura de mérito} & \textbf{Unidad} & \textbf{Cómo se mide} \\
    \midrule
    Tiempo de ciclo del cambio & horas & Del registro del cambio en el
      repositorio a su aplicación efectiva en producción. \\
    Frecuencia de aplicación & cambios/semana & Recuento de aplicaciones
      correctas por entorno. \\
    Tasa de fallo del cambio & \% & Aplicaciones que exigen reversión o
      corrección inmediata sobre el total. \\
    Tiempo de restablecimiento & minutos & De la detección del fallo a la
      recuperación del servicio. \\
    Cobertura de codificación & \% & Recursos gestionados por el motor sobre
      el total de recursos existentes en el proveedor. \\
    Tasa de deriva & recursos desviados / 1000 · semana & Diferencias
      detectadas en las ejecuciones de sólo lectura. \\
    Tiempo de reconstrucción & horas & De un entorno vacío a un entorno
      equivalente al de producción, partiendo sólo del repositorio. \\
    Cumplimiento de política & \% & Planes que superan la evaluación de
      políticas sin excepción manual. \\
    \bottomrule
  \end{tabularx}
\end{table}

El tiempo de reconstrucción es la figura más reveladora de todas, porque
verifica de un solo golpe las cuatro propiedades exigidas en la sección 1. Una
organización que no puede reconstruir un entorno completo desde el repositorio
no tiene un sistema de IaC completo, con independencia de cuántos ficheros
declarativos posea.

% ==========================================================================
\section{Modelo de objetos y procesos (OPM)}
% ==========================================================================

\subsection*{01. Por qué OPM para este sistema}

La metodología de objetos y procesos, normalizada como ISO/PAS 19450:2015 y
posteriormente como ISO 19450:2024
\parencite{iso19450:2015,iso19450:2024,dori2002,dori2016}, presenta dos rasgos que la
hacen adecuada aquí. El primero es que usa un único tipo de diagrama para
representar a la vez estructura y comportamiento, lo que evita el problema
habitual de mantener sincronizados un diagrama de bloques y un diagrama de
actividad. El segundo es su carácter bimodal: cada diagrama tiene una
traducción automática a lenguaje natural controlado (OPL), lo que permite
verificar el modelo leyéndolo.

Para un sistema de IaC, cuya esencia es un proceso que transforma el estado de
un objeto (la infraestructura desplegada) usando otros objetos como
instrumentos, la correspondencia con el vocabulario de OPM es casi directa.

\subsection*{02. Entidades del diagrama de sistema}

El Cuadro~\ref{tab:entidades} enumera las entidades del nivel superior.

\begin{table}[ht]
  \centering
  \footnotesize
  \caption{Entidades del diagrama de sistema (SD) del modelo OPM.}
  \label{tab:entidades}
  \begin{tabularx}{\textwidth}{@{}>{\raggedright\arraybackslash}p{0.27\textwidth}
    >{\raggedright\arraybackslash}p{0.13\textwidth}
    >{\raggedright\arraybackslash}X@{}}
    \toprule
    \textbf{Entidad} & \textbf{Tipo} & \textbf{Papel en el modelo} \\
    \midrule
    Aprovisionamiento de Infraestructura & Proceso & Proceso principal del
      sistema; se amplía en detalle en el diagrama SD1. \\
    Equipo de Plataforma & Objeto (físico) & Agente humano que maneja el
      proceso. \\
    Especificación de Infraestructura & Objeto & El código; instrumento del
      proceso. Se especializa en módulo reutilizable y configuración de
      entorno, y exhibe el atributo Versión. \\
    Repositorio Versionado & Objeto & Todo del que la especificación es
      parte. \\
    Motor de IaC & Objeto & Instrumento que ejecuta la conciliación. \\
    Registro de Estado & Objeto & Objeto afectado: el proceso lo lee y lo
      reescribe en cada aplicación. \\
    Proveedor de Nube & Objeto (físico) & Instrumento: expone la interfaz de
      programación. \\
    Credenciales & Objeto & Instrumento con el que el motor se autentica. \\
    Conjunto de Políticas & Objeto & Instrumento de la validación. \\
    Infraestructura Desplegada & Objeto (físico) & Operando del sistema.
      Estados: inexistente, conforme, desviada, retirada. \\
    Plan de Cambios & Objeto & Resultado del proceso, consumido por la
      aplicación. \\
    Registro de Auditoría & Objeto & Resultado del proceso. \\
    \bottomrule
  \end{tabularx}
\end{table}

La notación se reparte en dos figuras para que ninguna quede saturada. La
Figura~\ref{fig:opmsd} recoge los enlaces procedimentales, es decir, los que
unen el proceso con los objetos que lo habilitan y con los que transforma:
círculo hueco para instrumento, círculo relleno para agente, flecha doble
para efecto y flecha simple para resultado. Los enlaces estructurales se
tratan aparte, en el apartado 04. En ambas figuras, el rectángulo denota
objeto, la elipse denota proceso y el rectángulo de esquinas redondeadas
denota estado.

\begin{figure}[H]
  \centering
  \begin{tikzpicture}[scale=0.9, transform shape]
    % ---- proceso principal
    \node[opmproc, minimum width=40mm, minimum height=15mm] (P) at (0,0)
      {Aprovisionamiento\\de Infraestructura};

    % ---- agente humano
    \node[opmobj, fill=green!18] (eq) at (0,2.7) {Equipo de\\Plataforma};
    \draw[agente] (eq) -- (P);

    % ---- instrumentos
    \node[opmobj] (spec) at (-5.4,1.6)  {Especificación de\\Infraestructura};
    \node[opmobj] (motor) at (-5.4,0)   {Motor de IaC};
    \node[opmobj] (pol) at (-5.4,-1.6)  {Conjunto de\\Políticas};
    \node[opmobj] (prov) at (5.4,1.6)   {Proveedor\\de Nube};
    \node[opmobj] (cred) at (5.4,0)     {Credenciales};
    \foreach \nd in {spec,motor,pol,prov,cred} \draw[instr] (\nd) -- (P);

    % ---- objeto afectado
    \node[opmobj] (edo) at (-5.4,-3.2) {Registro de\\Estado};
    \draw[efecto] (edo) -- (P);

    % ---- resultados
    \node[opmobj] (plan) at (5.4,-1.6) {Plan de Cambios};
    \node[opmobj] (aud) at (5.4,-3.2)  {Registro de\\Auditoría};
    \draw[flecha] (P) -- (plan);
    \draw[flecha] (P) -- (aud);

    % ---- operando con sus cuatro estados
    \node[draw, very thick, fill=green!8, minimum width=76mm,
      minimum height=17mm] (infra) at (0,-4.6) {};
    \node[anchor=north, font=\footnotesize\bfseries]
      at ([yshift=-2mm]infra.north) {Infraestructura Desplegada};
    \foreach \dx/\st/\fl in {-27/inexistente/white, -9/conforme/{green!25},
      9/desviada/white, 27/retirada/white}
      \node[opmstate, fill=\fl]
        at ([xshift=\dx mm, yshift=-4.5mm]infra.center) {\st};
    \draw[efecto] (P) -- (infra);

    % ---- leyenda de notación
    \node[draw=gray!60, fill=gray!5, rounded corners=2pt, inner sep=4pt,
      font=\scriptsize, align=left, anchor=north, text width=0.95\textwidth]
      at ([yshift=-4mm]infra.south)
      {\textbf{Notación:} rectángulo = objeto · elipse = proceso ·
       rectángulo redondeado = estado · línea acabada en círculo hueco =
       instrumento · línea acabada en círculo relleno = agente · flecha
       doble = efecto · flecha simple = resultado};
  \end{tikzpicture}
  \caption{Diagrama de sistema (SD) del sistema de infraestructura como
  código, en notación OPM. Elaboración propia siguiendo
  \textcite{dori2016} y \textcite{iso19450:2024}.}
  \label{fig:opmsd}
\end{figure}

\subsection*{03. El lenguaje OPL del diagrama de sistema}

La traducción del diagrama al lenguaje de objetos y procesos permite
verificarlo por lectura. Las palabras reservadas van en versalitas.

\begin{quote}\small
\textsc{Equipo de Plataforma} es físico.\\
\textsc{Infraestructura Desplegada} es física.\\
\textsc{Proveedor de Nube} es físico y externo.\\
\textsc{Repositorio Versionado} consta de \textsc{Especificación de
Infraestructura}.\\
\textsc{Especificación de Infraestructura} puede ser \textsc{Módulo
Reutilizable} o \textsc{Configuración de Entorno}.\\
\textsc{Especificación de Infraestructura} exhibe \textsc{Versión}.\\
\textsc{Infraestructura Desplegada} puede ser inexistente, conforme, desviada
o retirada.\\
\textsc{Aprovisionamiento de Infraestructura} requiere \textsc{Especificación
de Infraestructura}, \textsc{Motor de IaC}, \textsc{Conjunto de Políticas},
\textsc{Proveedor de Nube} y \textsc{Credenciales}.\\
\textsc{Equipo de Plataforma} maneja \textsc{Aprovisionamiento de
Infraestructura}.\\
\textsc{Aprovisionamiento de Infraestructura} afecta a \textsc{Registro de
Estado}.\\
\textsc{Aprovisionamiento de Infraestructura} cambia \textsc{Infraestructura
Desplegada} de inexistente a conforme.\\
\textsc{Aprovisionamiento de Infraestructura} produce \textsc{Plan de Cambios}
y \textsc{Registro de Auditoría}.
\end{quote}

Leído en voz alta, el conjunto de frases es una descripción funcional completa
del sistema, y cualquier error de modelado se manifiesta como una frase falsa.

\subsection*{04. Los enlaces estructurales}

Además de participar en el proceso, los objetos del modelo mantienen entre sí
relaciones que no dependen de ninguna ejecución. La
Figura~\ref{fig:opmestr} recoge las tres que afectan a la \emph{Especificación
de Infraestructura}: triángulo relleno para la agregación que la vincula con
el repositorio que la contiene, triángulo hueco para la generalización en sus
dos especializaciones y triángulo doble para la exhibición del atributo que la
versiona. Son las relaciones que sostienen las frases \emph{consta de},
\emph{puede ser} y \emph{exhibe} del OPL anterior.

\begin{figure}[H]
  \centering
  \begin{tikzpicture}[scale=0.95, transform shape]
    \node[opmobj, text width=34mm] (spec) at (0,0)
      {Especificación de\\Infraestructura};
    % ---- agregación
    \node[opmobj, text width=30mm] (repo) at (0,2.1)
      {Repositorio\\Versionado};
    \draw[very thick] (repo.south) -- (spec.north);
    \coordinate (ca) at ($(repo.south)!0.5!(spec.north)$);
    \opmagg{0}{ca}
    \node[etiq, right=2mm] at (ca) {agregación: \emph{consta de}};
    % ---- exhibición del atributo
    \node[opmobj, text width=22mm] (ver) at (5.4,0) {Versión};
    \draw[very thick] (spec.east) -- (ver.west);
    \coordinate (ce) at ($(spec.east)!0.42!(ver.west)$);
    \opmexh{90}{ce}
    \node[etiq, above=5mm] at ($(spec.east)!0.6!(ver.west)$)
      {exhibición: \emph{exhibe}};
    % ---- generalización
    \node[opmobj, text width=30mm] (mod) at (-3.1,-2.6)
      {Módulo\\Reutilizable};
    \node[opmobj, text width=30mm] (cfg) at (3.1,-2.6)
      {Configuración\\de Entorno};
    \coordinate (cg) at (0,-1.25);
    \draw[very thick] (spec.south) -- (0,-1.6);
    \draw[very thick] (mod.north) |- (0,-1.6) -| (cfg.north);
    \opmgen{0}{cg}
    \node[etiq, right=2mm] at (cg) {generalización: \emph{puede ser}};
  \end{tikzpicture}
  \caption{Enlaces estructurales del objeto \emph{Especificación de
  Infraestructura}: agregación con el repositorio que la contiene,
  generalización en sus dos especializaciones y exhibición del atributo que
  la versiona. Elaboración propia.}
  \label{fig:opmestr}
\end{figure}

\subsection*{05. Ampliación en detalle del proceso principal}

El diagrama SD1 amplía el proceso principal en siete subprocesos, ordenados
verticalmente según la convención de OPM, en la que la posición vertical
denota precedencia temporal. Los objetos que sólo intervienen en el interior
del proceso se dibujan dentro de la elipse; el \emph{Equipo de Plataforma},
que pertenece al entorno, queda fuera. La Figura~\ref{fig:opmsd1} lo recoge.

\begin{figure}[H]
  \centering
  \begin{tikzpicture}[scale=0.82, transform shape, x=1mm, y=1mm]
    \tikzset{sp/.style={opmproc, minimum width=34mm, minimum height=10mm,
               text width=},
             ob/.style={opmobj, text width=26mm, minimum height=9mm,
               font=\scriptsize\bfseries}}
    % ---- marco elíptico del proceso ampliado
    \node[draw, very thick, ellipse, fill=blue!4,
      minimum width=130mm, minimum height=146mm] (marco) at (0,3) {};
    \node[font=\footnotesize\bfseries] at (0,62)
      {Aprovisionamiento de Infraestructura};

    % ---- los siete subprocesos, de arriba abajo
    \node[sp] (c)  at (0, 44) {Codificación};
    \node[sp] (v)  at (0, 30) {Validación};
    \node[sp] (pl) at (0, 16) {Planificación};
    \node[sp] (au) at (0,  2) {Autorización};
    \node[sp] (ap) at (0,-12) {Aplicación};
    \node[sp] (co) at (0,-26) {Conciliación};
    \node[sp] (re) at (0,-40) {Retirada};
    \foreach \a/\b in {c/v,v/pl,pl/au,au/ap,ap/co,co/re}
      \draw[flecha, gray!65] (\a) -- (\b);

    % ---- objetos internos
    \node[ob] (spec) at (-36, 30)
      {Especificación\\ \scriptsize\mdseries borrador $\to$ validada};
    \node[ob] (pol)  at ( 36, 30) {Conjunto de\\Políticas};
    \node[ob] (edo)  at (-36, 16) {Registro\\de Estado};
    \node[ob] (plan) at ( 36, 16)
      {Plan de Cambios\\ \scriptsize\mdseries pendiente $\to$ aprobado};
    \node[ob] (cred) at ( 36, -6) {Credenciales};
    \node[ob] (prov) at ( 36,-18) {Proveedor\\de Nube};
    \node[ob] (aud)  at ( 36,-31) {Registro de\\Auditoría};
    \node[draw, very thick, fill=green!8, minimum width=32mm,
      minimum height=14mm] (infra) at (-36,-26) {};
    \node[anchor=north, font=\scriptsize\bfseries, align=center]
      at ([yshift=-1mm]infra.north) {Infraestructura\\Desplegada};
    \node[opmstate, font=\tiny] at ([xshift=-7mm,yshift=-4.5mm]infra.center)
      {desviada};
    \node[opmstate, font=\tiny, fill=green!25]
      at ([xshift=7mm,yshift=-4.5mm]infra.center) {conforme};

    % ---- agente humano, fuera del proceso ampliado
    \node[ob, fill=green!18, text width=24mm] (eq) at (-70, 50)
      {Equipo de\\Plataforma};
    \draw[agente] (eq) to[out=0,in=180] (c.west);
    \draw[agente] (eq.south) |- (au.west);

    % ---- enlaces de los subprocesos
    \draw[flecha] (c.west) to[out=180,in=90] (spec.north);
    \draw[efecto] (spec) -- (v);
    \draw[instr]  (pol)  -- (v);
    \draw[instr]  (edo)  -- (pl);
    \draw[flecha] (pl)   -- (plan);
    \draw[instr]  (plan.south) to[out=-90,in=0] (au.east);
    \draw[instr]  (cred) -- (ap);
    \draw[instr]  (prov) -- (ap);
    \draw[efecto] (ap.west) to[out=180,in=90] (infra.north);
    \draw[efecto] (co) -- (infra);
    \draw[flecha] (co) -- (aud);
    \draw[efecto] (re.west) to[out=180,in=-90] (infra.south);
  \end{tikzpicture}
  \caption{Diagrama SD1: ampliación en detalle del proceso de
  aprovisionamiento. La posición vertical de los subprocesos denota
  precedencia temporal. \emph{Aplicación} y \emph{Conciliación} actúan sobre
  el mismo objeto desde estados de partida distintos, que es la formulación
  precisa de la diferencia entre desplegar y corregir la deriva. El
  \emph{Equipo de Plataforma} queda fuera del proceso ampliado por ser un
  objeto del entorno. Elaboración propia.}
  \label{fig:opmsd1}
\end{figure}

\subsection*{06. El lenguaje OPL de la ampliación}

\begin{quote}\small
\textsc{Aprovisionamiento de Infraestructura} consta de
\textsc{Codificación}, \textsc{Validación}, \textsc{Planificación},
\textsc{Autorización}, \textsc{Aplicación}, \textsc{Conciliación} y
\textsc{Retirada}.\\
\textsc{Equipo de Plataforma} maneja \textsc{Codificación} y
\textsc{Autorización}.\\
\textsc{Codificación} produce \textsc{Especificación de Infraestructura}.\\
\textsc{Validación} requiere \textsc{Conjunto de Políticas}.\\
\textsc{Validación} cambia \textsc{Especificación de Infraestructura} de
borrador a validada.\\
\textsc{Planificación} requiere \textsc{Registro de Estado} y
\textsc{Especificación de Infraestructura}.\\
\textsc{Planificación} produce \textsc{Plan de Cambios}.\\
\textsc{Autorización} cambia \textsc{Plan de Cambios} de pendiente a
aprobado.\\
\textsc{Aplicación} requiere \textsc{Plan de Cambios}, \textsc{Credenciales}
y \textsc{Proveedor de Nube}.\\
\textsc{Aplicación} cambia \textsc{Infraestructura Desplegada} de inexistente
a conforme.\\
\textsc{Aplicación} afecta a \textsc{Registro de Estado}.\\
\textsc{Conciliación} cambia \textsc{Infraestructura Desplegada} de desviada
a conforme.\\
\textsc{Conciliación} produce \textsc{Registro de Auditoría}.\\
\textsc{Retirada} cambia \textsc{Infraestructura Desplegada} de conforme a
retirada.
\end{quote}

Merece la pena señalar lo que el modelo hace explícito y que las descripciones
en prosa suelen dejar implícito. Primero, que \textsc{Registro de Estado} es
el único objeto que aparece a la vez como instrumento de la planificación y
como objeto afectado por la aplicación, lo que identifica exactamente por qué
necesita bloqueo de concurrencia. Segundo, que \textsc{Conciliación} y
\textsc{Aplicación} cambian el mismo objeto desde estados de partida
distintos, que es la formulación precisa de la diferencia entre desplegar y
conciliar. Y tercero, que el equipo humano interviene sólo en dos de los siete
subprocesos, lo que delimita el alcance real de la automatización.

% ==========================================================================
\section{Conclusiones}
% ==========================================================================

\begin{enumerate}
  \item El funcionamiento de un sistema de infraestructura como código se
    reduce a un mecanismo único: la conciliación repetida entre un estado
    deseado, expresado en código versionado, y un estado real, observado a
    través de la interfaz de un proveedor. Todo lo demás, herramientas,
    formatos y flujos de trabajo, son variaciones sobre ese mecanismo.
  \item La pieza que hace posible el mecanismo es el estado registrado. No es
    un detalle de implementación sino el elemento que permite reconocer la
    propiedad de los recursos, detectar destrucciones y encadenar dependencias,
    y por eso concentra los requisitos de cifrado, versionado y bloqueo.
  \item Un sistema completo comprende siete capas, del sustrato del proveedor
    a la interfaz de autoservicio, y seis módulos transversales. Las capas
    intermedias compiten por el mismo territorio, de modo que la primera
    decisión de arquitectura consiste en asignar de forma explícita la
    propiedad de cada característica a una sola capa.
  \item Los módulos transversales, y en particular la política como código y
    las pruebas, son los que distinguen un sistema que automatiza de uno que
    además garantiza. Su ausencia no impide desplegar: hace que los errores se
    desplieguen igual de rápido que los aciertos.
  \item El tiempo de reconstrucción de un entorno completo desde el
    repositorio es la figura de mérito que mejor resume el estado de una
    implantación, porque verifica de una vez las cuatro propiedades que
    definen al sistema.
  \item El modelado en OPM aporta algo que la descripción en prosa no aporta:
    obliga a declarar el tipo de cada relación. Al hacerlo aparecen con
    claridad la doble condición del registro de estado, la diferencia formal
    entre aplicar y conciliar, y el alcance exacto de la intervención humana,
    que en este modelo se limita a dos de los siete subprocesos.
\end{enumerate}

\printbibliography[title={Referencias bibliográficas}]

\end{document}
